\section{Generative evaluation workload}
All evaluations in this report are generative: the model writes an answer
through its chat template and the saved answer is graded. Experiment design,
model selection, dataset revisions, and generation settings are defined in
\texttt{experiments/qwen3\_mmlu/run.py}; the paper does not maintain a second
configuration. Activation zeroing operates at the selected experts' MLP
intermediates, the residual-stream sum, or the MoE output. Weight zeroing
operates on checkpoint parameters. Their accuracy results require comparisons
against each benchmark's own baseline and repeat.

Tables~\ref{tab:input-lengths} and~\ref{tab:output-lengths} describe the
completed baseline workloads. Each column is a benchmark, each row is an
experiment arm, and each mean weights its questions equally. The repeat uses
the same protocol and captures run-to-run variation. Benchmark sizes appear in
Table~\ref{tab:sample-counts}. We do not pool benchmarks with different sizes.

Input length counts the saved initial prompt, including chat-template tokens.
Output length counts the saved model-written text, retokenized with the
experiment's pinned tokenizer and no additional special tokens. For MMLU STEM
and GSM8K, reasoning and any forced-answer continuation are encoded separately
and their lengths summed. The injected forced-answer suffix is excluded.
Invisible stop tokens are excluded as well. Thus these are text lengths rather
than server-reported completion usage. Input length does not include a second
prefill when an answer is forced. The analysis is offline and requires no new
model generation.

\begin{table}[ht]
\centering
\caption{Mean initial input length in tokens.}
\label{tab:input-lengths}
% Generated by scripts/paper/token_lengths.py; do not edit.
\resizebox{\textwidth}{!}{%
\begin{tabular}{lrrrrrrrrrrr}
\toprule
Experiment & MMLU STEM & GSM8K & MATH500 & IFEval & MMLU-Pro & GPQA & HumanEval+ & MBPP+ & RULER 8k & RULER 16k & RULER 32k \\
\midrule
baseline & 135.0 & 113.3 & 100.9 & 54.4 & 217.8 & 239.8 & 164.2 & 81.5 & 7742.8 & 15911.2 & 32082.6 \\
baseline\_repeat & 135.0 & 113.3 & 100.9 & 54.4 & 217.8 & 239.8 & 164.2 & 81.5 & 7742.8 & 15911.2 & 32082.6 \\
\bottomrule
\end{tabular}}

\end{table}
\begin{table}[ht]
\centering
\caption{Mean model-written output text length in tokens.}
\label{tab:output-lengths}
% Generated by scripts/paper/token_lengths.py; do not edit.
\resizebox{\textwidth}{!}{%
\begin{tabular}{lrrrrrrrrrrr}
\toprule
Experiment & MMLU STEM & GSM8K & MATH500 & IFEval & MMLU-Pro & GPQA & HumanEval+ & MBPP+ & RULER 8k & RULER 16k & RULER 32k \\
\midrule
baseline & 205.1 & 112.1 & 1241.0 & 439.8 & 807.6 & 1817.0 & 181.1 & 141.3 & 29.0 & 28.7 & 28.5 \\
baseline\_repeat & 198.5 & 111.5 & 1266.3 & 425.3 & 807.6 & 1852.5 & 200.9 & 141.3 & 29.0 & 28.9 & 28.6 \\
\bottomrule
\end{tabular}}

\end{table}

\section{Interpretation and limitations}
Input length describes prefill workload, while output length describes the
amount of generated text. Neither is a latency measurement: kernel performance,
batching, prefix caching, and repeated prefills are not measured here.
Different task instructions and generation caps also affect output length.
A response ending at a generation cap remains in the mean; the table describes
observed responses, not an estimate of uncapped answer length. In particular,
RULER columns use their configured mixture of subtasks and answer budgets.
